\begin{table}[!htbp]
\centering
\caption{Block-permutation null on exchangeable and recorded data (family-wise rejection rate at $\alpha = 0.05$, 199 permutations, 300 trials per cell, $d = 20$, windows of 400 rows). A rate above the nominal $5\%$ means the permuted null is too optimistic; block size $1$ is the textbook permutation and larger blocks keep adjacent rows together. Pooled synthetic rows are i.i.d. by construction, so a rate that moves with block size there is Monte-Carlo noise; the recorded streams keep their order, so those rows are the measurement.}
\label{tab:block-null}
\footnotesize
\setlength{\tabcolsep}{4pt}
\resizebox{\columnwidth}{!}{%
\begin{tabular}{lrrrr}
\toprule
Pool & block $1$ & block $10$ & block $25$ & block $50$ \\
\midrule
\multicolumn{5}{l}{\emph{nothing injected (nominal $5\%$)}} \\
Gaussian (i.i.d.) & 3.3 & 3.3 & 3.7 & 3.3 \\
Multivariate $t_3$ (i.i.d.) & 3.3 & 1.7 & 4.0 & 2.7 \\
CICIDS Tuesday, recorded & 88.7 & 69.3 & 50.0 & 41.3 \\
NSL-KDD, random order & 1.7 & 2.3 & 4.3 & 2.7 \\
\midrule
\multicolumn{5}{l}{\emph{one feature shifted by $+0.5\sigma$}} \\
Gaussian (i.i.d.) & 98.7 & 99.3 & 99.0 & 98.3 \\
Multivariate $t_3$ (i.i.d.) & 93.7 & 92.3 & 94.0 & 93.7 \\
CICIDS Tuesday, recorded & 100.0 & 100.0 & 100.0 & 100.0 \\
NSL-KDD, random order & 100.0 & 100.0 & 100.0 & 100.0 \\
\bottomrule
\end{tabular}}
\end{table}
